%!TEX root = hems_proj.tex
%%%%%%%%%%%%%%%%%%%%%%%%%%%%%%%%%%%%%%%%%%%%%%%%%%%%%%%%%%%%%%%%%%%%%%%
\chapter{Módszertan és Rendszerarchitektúra}\label{ch:MODSZERTAN}
%%%%%%%%%%%%%%%%%%%%%%%%%%%%%%%%%%%%%%%%%%%%%%%%%%%%%%%%%%%%%%%%%%%%%%%

\begin{osszefoglal}
Ez a fejezet a projekt jelenlegi implementációs állapotának megfelelően formalizálja a teljes pipeline-t: adatforrások, profilépítés, optimalizálási ciklus és kiértékelési lépések. A cél, hogy az eredmények ne csak reprodukálhatók, hanem módszertanilag auditálhatók is legyenek.
\end{osszefoglal}

\section{Rendszerblokkok és felelősségek}

A megoldás moduláris felépítésű, az egyes komponensek egyértelmű felelősségi határokkal működnek. A \texttt{hems\_problem.py} modul tartalmazza a döntési teret, a célfüggvényt és a korlátkezelést. Az \texttt{algorithms.py} fájlban kapott helyet a GA, PSO, GWO és a hibrid GA-PSO optimalizáló. A \texttt{data\_generator.py} a szintetikus load/PV profilok generálásáért és a CSV validációért felelős, míg a \texttt{price\_loader.py} az ENTSO-E árak betöltését és a napi árprofil-képzést végzi. A teljes benchmark-vezérlés, az aggregáció és a riportolás a \texttt{benchmark\_real\_vs\_synthetic.py} modulban van összefogva. A moduláris szerkezet célja, hogy egy-egy részegység -- például az árkezelés vagy a célfüggvény -- változtatása minimális oldalkövetkezménnyel járjon.

\section{Adatforrások és előfeldolgozás}

A benchmark két fő külső adatkört használ: valós inverteres teljesítmény-idősorokat (5 perces mintavételből képzett órás átlag), valamint ENTSO-E day-ahead árakat (EUR/MWh), amelyek RON/kWh formára konvertálódnak. Az áradatok feldolgozása során területi szűrés (RO zóna), időbélyeg-normalizálás és órás újramintavételezés történik. Hiányos napi árprofil esetén fallback TOU ár kerül használatra, ami biztosítja, hogy részlegesen hiányos forrás mellett is végigfutható maradjon a pipeline, és transzparensen jelzi, mely napoknál történt visszaesés az alapértelmezett árra.

\section{Profilépítés: A 2×2-es mátrix}

A benchmark minden kiválasztott napra négy profilkombinációt képez a termelési és fogyasztási adatokból \cite{Lazzari2023}. Az első kombináció szintetikus PV-t és szintetikus loadot párosít, ez a laboratóriumi ideális eset. A második valós inverter PV-t és szintetikus loadot használ, így elkülöníti a termelési oldal bizonytalanságának hatását. A harmadik kombináció szintetikus PV mellé valós UCI háztartási fogyasztást rendel, a negyedik pedig mindkét oldalon valós adatokat alkalmaz -- ez a teljes valóság forgatókönyve. A döntés mögötti módszertani cél a bemeneti realizmus és a különböző zajforrások hatásának elkülönítése az algoritmusok teljesítményétől.

\section{Optimalizálási ciklus és futásszám}

Jelölje $D$ a napok számát, $S$ a dataset-változatok számát, $R$ a replikák számát, $A$ az algoritmusok számát, $K$ pedig a futások számát. Ekkor a teljes futásszám:
\begin{equation}
N_{eval} = D \cdot S \cdot R \cdot A \cdot K
\end{equation}
A jelenlegi mátrix-konfigurációban $D=16$, $S=4$, $R=1$, $A=5$, $K=2$, ezért összesen \textbf{640 algoritmus-futtatás} történik. Ez a szám elegendő a rangsorok megbízható megállapításához, ugyanakkor végleges statisztikai következtetésekhez bővebb futásszám ajánlott.

\section{Kiértékelési mutatók}

A nyers futásszintű adatokból algoritmusonként és adathalmazonként több aggregált mutató képződik: az átlagköltség és szórása, a legjobb költség, az átlagos futási idő, az átlagos relatív rés a profilonkénti legjobbhoz képest, valamint a győzelmek száma és aránya. A gap egy profile-szintű normalizált mutató:
\begin{equation}
gap_{i,a}[\%] = \frac{c^{best}_{i,a} - c^{best}_{i,*}}{|c^{best}_{i,*}|+\varepsilon} \cdot 100
\end{equation}
ahol $c^{best}_{i,*}$ az adott profil globális legjobb költsége. Ez a mutató lehetővé teszi, hogy az algoritmusok teljesítményét ne csak abszolút értékben, hanem relatív versenyképességük alapján is össze lehessen hasonlítani.